\scanpage{72}
% archive foliation: 71

Soit \deleted{$\unsure{\{\mathcal C\}}$ \unsure{une}}
$\mathcal O$ la catégorie des $A$-algèbres \unsure{apparentées}
\deleted{\unsure{centrales}} à une $K$-Algèbre déterminée,
\deleted{et \unsure{pour tout}}
\deleted{\unsure{Pour tout} $O\in\mathcal O$, \unsure{supposons qu'il existe une fonction}}
(où $K$ l'anneau total des
\deleted{\unsure{Soit $A$ rég. loc.}}
fractions de $A$, \inserted{d'un $A\subset K$}
et supposons $K$ $\frac12$ loc.

Donc \deleted{\unsure{tt}} pour $O\in\mathcal O$, $O'\in\mathcal O$,
tt bimodule $L$ sur $O,O'$ loc. libre de rang $1$
\inserted{de chaque côté} loc. sur $S=\operatorname{Spec}A$ est libre
\inserted{\unsure{génériquement}}
sur $O_K\unsure{\simeq}E$ et sur $O'_K$, d'où on conclut que
\inserted{\uuline{si tout autom de $E$ est intérieur}}
a) $O'_K\simeq E$ \deleted{\unsure{donc $L$}}
\inserted{donc $O'$ se plonge dans $E$}
et b) tt $L$ ci-dessus se plonge dans $E$ comme un $(O,O')$-biidéal
\inserted{avec la structure \unsure{d'alg.}}

\margin{\unsure{$E$ est donc pris au départ pour un quelconque membre de $\mathcal O$}.}

La \deleted{\unsure{sur}} $2$-catégorie
\inserted{pleine de $\unsure{\underline B}\ \unsure{(\mathrm{Cat})}$}
dont les objets sont \deleted{\unsure{les $O$, les}}
\inserted{les cat. de}
Deuring de type $\mathcal O$, \deleted{et}
qui ont \unsure{gén. strict} : \unsure{alors} $\uline{\underline B_1}$
dont les objets sont les objets de $\mathcal O$,
\[
\uline{\operatorname{Hom}}(O,O')
=\inserted{\text{catégorie des}}\text{bimodules sur }O,O'
\text{ qui sont loc. }\inserted{\text{libres de}}\deleted{\text{rang}}\text{ rang }1
\]
de chaque côté loc. sur $S$, et \unsure{l'accouplement}
\[
\deleted{\unsure{{}_OM_{O'},{}_{O'}L_{O''}\longrightarrow M\otimes_{O'}L}}
\]
\[
{}_OL_{O'},{}_{O'}M_{O''}\longrightarrow L\otimes_{O'}M,
\qquad \uline{\underline B_2}
\]
est aussi équivalente au \uline{$2$-groupoïde de Brandt}
dont les objets sont les sous-anneaux de $E$
\deleted{\unsure{algèbres}}
$\in\mathcal O$,
\[
\uline{\operatorname{Hom}}(O,O')=
\deleted{\unsure{catégorie des}}\inserted{\text{ensemble des}}
(O,O')\text{-biidéaux }L
\]
de $E$ \unsure{lesquels} des deux côtés loc. lib. de rg $1$ loc. sur $S$
\[
\operatorname{Hom}(L,M)=\text{ens des }\lambda\in
\unsure{E^{Z*}|E^\sharp}\text{ tels que }\lambda L=M
\]
\inserted{\unsure{centre de $E$}}
\inserted{si $L,M$ sont les $(O,O')$ biidéaux ``\unsure{spéciaux}''}.

On notera que les $2$-foncteurs
\[
\underline B_2\longrightarrow\underline B_1\longrightarrow
\underline B\equiv\underline B_0
\]
n'ont pas de quasi-inverses ``naturels'' !
